\section{Research Method and Evaluation}

Three questions guide evaluation: (RQ1) How much does temporal pose modeling improve on geometry rules? (RQ2) What changes in F1, false alerts/hour, and delay result from quality-aware fusion and the alert state machine? (RQ3) What throughput, latency, and bandwidth does the edge--cloud split achieve?

\begin{table}[H]
\caption{Controlled comparison and purpose.}
\label{tab:ablation}
\centering
\footnotesize
\begin{tabularx}{\columnwidth}{@{}p{0.10\columnwidth}X@{}}
\toprule
ID & System variant \\
\midrule
B0 & Box/keypoint geometry rule \\
B1 & TCNTE-style skeleton temporal baseline \\
B2 & Skeleton and geometry branches, fixed fusion \\
B3 & B2 plus learned pose-quality gate \\
B4 & B3 plus multi-stage alert state machine \\
\bottomrule
\end{tabularx}
\end{table}

UP-Fall uses participant-grouped cross-validation. Cross-dataset evaluation on GMDCSA-24 will examine generalization beyond the primary benchmark. Fall metrics are precision, recall, F1, specificity, event recall, false alerts/hour, and detection delay. Robustness will be evaluated under controlled visual degradation, including reduced illumination, missing or occluded keypoints, and frame loss. Transformations and random seeds will be recorded.

CPU and Hailo-10H comparisons will use the same input size and report throughput, edge latency, resource use, and temperature. Network and cloud measurements will cover payload size and end-to-end latency. Replayed fall events will exercise alerting, storage, dashboard updates, acknowledgement, and reconnection. The final evaluation will include continuous camera operation, cloud connectivity, persisted events, B0--B4 model comparisons, robustness results, and reproducible run instructions.
